\documentclass[11pt]{article}




\usepackage[utf8]{inputenc}
\usepackage[T1]{fontenc}
\usepackage[margin=1in]{geometry}
\usepackage{amsmath,amssymb,amsthm}
\usepackage{booktabs}
\usepackage{graphicx}
\usepackage{microtype}
\usepackage[hidelinks]{hyperref}
\usepackage{float}


\newtheorem{lemma}{Lemma}

\title{\textbf{How Much of Truth Fits on a Single Axis?}\\[2pt]
\large Knowledge-Dependent Dimensionality, Polarity Recovery,\\
and the Anatomy of a Truth Direction in Qwen2.5-1.5B}

\author{Francesco Karim Vicidomini\\
\small Independent researcher\\
\small Reference code: \texttt{truth\_probe.py}, \texttt{canvas.py}, \texttt{anatomy.py} and companion tools}

\date{July 2026\thanks{This version consolidates two earlier
manuscripts, ``How Much of Truth Fits on a Single Axis?'' (June 28, 2026) and
the working note ``Anatomy of the Truth Axis in Qwen2.5-1.5B'', into a single
revised document.}}

\begin{document}
\maketitle

\begin{abstract}
Reading whether a language model represents a statement as true or false
usually relies on trained linear probes or on sparse autoencoders. We study a
minimal alternative: a single direction obtained, with no training, from the
singular value decomposition (SVD) of hidden-state differences over
true/false minimal pairs; the direction is identified without labels up to
one global sign (Lemma~\ref{lem:labels}). On Qwen2.5-1.5B, under a
conservative held-out protocol, we report four findings. First, the axis
carries a real truth signal (held-out AUC $0.938$), yet a full linear probe
reaches $0.989$: one direction does not exhaust the representation. Second,
the gap between axis and probe grows as the model's knowledge of the facts
decreases, and grows further when heterogeneous material forces one axis to
serve many truth frames: the dimensionality of the truth representation is
knowledge-dependent. Third, on mixed-polarity data the same decomposition
recovers the supervised polarity direction of B\"urger et al.\ (cosine
$0.959$) without seeing polarity labels. Fourth, attention builds the signal
in the middle layers and the feed-forward blocks erode it after the peak, a
claim supported causally by ablation. Seven falsification experiments
designed to reveal richer structure fail to overturn the one-dimensional
reading. Results are from a single small model and are intended to be
falsifiable.
\end{abstract}

\section{Introduction}

Large language models produce false statements with the same fluency as true
ones. Predicting when a model is about to produce a false statement, ideally
during generation and without a second judge model, is a central reliability
problem. Mechanistic interpretability has shown that many high-level
properties, including the truth value of a statement, are encoded as linear
directions in the residual stream~\cite{marks2024,zou2023}.

Two families of methods have been developed to read the truth value of a
statement from hidden states. The first is the \emph{linear probe}: a
classifier or difference-of-means axis whose projection separates true from
false~\cite{burns2023,marks2024,li2023}. The second is \emph{sparse dictionary
learning}: project the state onto an overcomplete dictionary, potentially
millions of atoms, to undo superposition and read monosemantic
features~\cite{cunningham2023,bricken2023}.

Each pays a price. A fixed axis reads a single linear functional of the state,
and prior work documents a concrete failure mode of such probes under
negation~\cite{burger2024}; a sparse autoencoder requires dedicated training
and a large matrix multiply per token. This paper takes a third stance. Rather
than learning a probe or expanding a dictionary, we read structure from
geometry alone, and ask a precise question: \emph{how much of the truth
structure does a single direction capture, and when does that capture fail?}
We extract all hidden states in one forward pass and let an axis emerge from
the SVD of paired true/false differences. Part of the apparent accuracy of any
such pipeline could come from selecting the best layer after the fact; we
therefore measure everything out of sample and quantify the selection effect
directly with a permutation null (Section~\ref{sec:protocol}). The answer is
modest but, we believe, clean: the axis reads truth cheaply when the fact is
known, fails predictably when it is not, and, blind to the label, recovers the
polarity structure that a prior supervised method constructed by hand.

A single positive result invites the obvious objection: is the
one-dimensionality real, or merely the easiest thing to see? The second half
of the paper responds with a systematic robustness analysis. We subject the
one-dimensional claim to seven falsification experiments, each designed to
expose a richer structure if one exists, and report that none overturns it. We
then go one level deeper and ask \emph{where} the axis comes from: decomposing
the residual stream into its attention and feed-forward (FFN) contributions,
we locate the components that build, carry, and erode the signal, and test the
resulting picture causally by ablation.

\paragraph{Contributions.}
\begin{enumerate}
\item A training-free truth probe at cost $O(d)$, where $d$ is the
residual-stream width, identified without labels up to one global sign
(Lemma~\ref{lem:labels}), with a held-out and permutation-tested evaluation
that corrects the optimism of in-sample axis selection
(Sec.~\ref{sec:signal}).
\item Evidence that the \emph{dimensionality of the truth representation is
knowledge-dependent}: near one-dimensional for known facts, diffuse for
unknown ones, quantified by the gap to a full linear probe
(Sec.~\ref{sec:dim}).
\item \emph{Recovery of the polarity direction} $t_P$ of B\"urger et
al.~\cite{burger2024} without polarity supervision: the SVD of mixed-polarity
differences aligns with the supervised $t_P$ at cosine $0.959$
(Sec.~\ref{sec:polarity}).
\item A \emph{robustness battery}: seven falsification experiments against the
one-dimensional reading (a second dimension, depth-global structure, dynamics,
layer departments, energy, phase, three-dimensional clustering), each tested
under held-out cross-validation and a permutation null, none of which
overturns the claim (Sec.~\ref{sec:robust}).
\item An \emph{anatomy} of the axis: an exact attention/FFN decomposition of
the residual stream with a built-in correctness gate, showing a stable
per-layer alternation in readability, a negative result inside the expanded
FFN space, and causal ablations establishing that attention builds the signal
in the middle layers while the FFN degrades it in the transition band
(Sec.~\ref{sec:anatomy}).
\item An interpretation of the residual overlap on unknown facts as the
observable signature of unresolved superposition, tied to a capacity argument
(Secs.~\ref{sec:discussion},~\ref{sec:why}).
\end{enumerate}

\section{Related work}

\paragraph{Linear truth directions.} \cite{burns2023} extract a knowledge
direction without labels by enforcing logical consistency across a statement
and its negation; the direction is found by optimization. \cite{marks2024}
show truth has emergent linear structure and that difference-of-means probes
generalize as well as more complex ones. \cite{li2023} intervene along
truthful directions at inference time; \cite{zou2023} frame representation
reading top-down. Our axis shares the linear-direction hypothesis, but it is
read from the geometry of paired differences by an SVD, with no optimization;
its precise relation to supervision is stated in Lemma~\ref{lem:labels}.

\paragraph{A two-dimensional truth subspace.} Closest to us, B\"urger et
al.~\cite{burger2024} show that truth is represented not in a single direction
but in a two-dimensional subspace spanned by a general truth direction $t_G$
and a polarity-sensitive direction $t_P \sim
\mathrm{XOR}(\text{is\_true},\text{is\_negated})$, and that probes trained on
affirmative statements fail on negated ones. We confirm this structure in
Qwen2.5-1.5B and show, as a methodological addition, that $t_P$ can be
recovered \emph{without} polarity supervision.

\paragraph{Cross-layer dynamics.} Recent monitoring methods model the
evolution of hidden states across depth rather than reading a single
layer~\cite{icr2025,hsad2025}. Our depth trajectories are descriptive and
consistent with this line.

\paragraph{Beyond single-direction probes on the residual stream.} All the
methods above, ours included, assume white-box access to the residual stream,
which is not always available and is not the only usable signal. Bar-Shalom et
al.~\cite{barshalom2025} treat the full activation tensor (layers by tokens)
as an image and train a vision-transformer detector on it, across several LLMs
simultaneously, with evaluation on a broad set of hallucination benchmarks.
Methods of this kind are more expressive than any single direction and are
trained for detection; our axis sits at the opposite end of the same
cost/expressivity curve, as the cheapest structural point, and the comparison
we draw is against that curve, not against any single method.

\paragraph{Superposition and dictionaries.} Superposition~\cite{elhage2022}
motivates decomposing the residual stream into overcomplete
dictionaries~\cite{cunningham2023,bricken2023}. This is the operation we
explicitly do \emph{not} perform. The bottom line, argued in
Secs.~\ref{sec:discussion} and~\ref{sec:why}, is that the two approaches
answer different questions and neither subsumes the other: a dictionary
answers ``which features are active'' at cost $O(d\,m)$ and after dedicated
training; the axis answers ``how much truth a single direction carries'' at
cost $O(d)$ with no training. When feature identity matters, dictionaries are
the right tool; for a cheap runtime signal on facts the model knows, the axis
suffices. We make no claim of superior accuracy over dictionary methods.

\section{Scope and evidence}
\label{sec:scope}

Everything in this paper concerns one small base model, Qwen2.5-1.5B, and
sentence-level factual statements. Results come in two evidential tiers, and
we keep them distinct throughout. \textbf{Tier 1} (Secs.~\ref{sec:results}
and~\ref{sec:robust}): the positive findings and the falsification battery,
evaluated held-out with permutation nulls. \textbf{Tier 2}
(Sec.~\ref{sec:anatomy}): the anatomy, a mechanistic follow-up with the same
protocol but smaller margins in places, single datasets for some entries, and
no falsification battery of its own; where a statement there is
interpretation rather than measurement, it is marked as such. Cross-model
generalization is not claimed here; what is known about it from subsequent
work is summarized, with numbers, in Section~\ref{sec:generalization}.

\section{Method}
\label{sec:method}

\subsection{Extraction}
\label{sec:extraction}

Informally: a transformer reads a sentence and, at every layer, updates an
internal vector summary of the text so far; we record that summary at the
position of the last token, where the sentence is complete, at every layer,
and study the geometry of these vectors.

Formally, we load Qwen2.5-1.5B frozen. The model has $28$ transformer blocks,
hence $29$ hidden-state levels (the embedding plus each block output), each a
vector of width $d = 1536$, the width of the residual stream. Per sentence we
run one forward pass and take the last-token state at every level. For $N$
sentences this yields $H \in \mathbb{R}^{N \times 29 \times d}$. Once an axis
is fixed, scoring a new state is a single dot product, hence the reading cost
of $O(d)$ per token quoted throughout.

\subsection{Data: minimal pairs}
\label{sec:data}

A \emph{minimal pair} is a pair of sentences with the same subject and the
same template that differ in exactly one content word, such that one sentence
is true and the other false: ``The Sun is a star.''\ vs.\ ``The Sun is a
planet.''\ Both sentences are affirmative. Within a pair, the difference of
hidden states cancels the shared topic and surface form and leaves the truth
signal.

Two sources are used. The \emph{curated pairs} are a built-in set of $45$
pairs of common facts (capitals, chemical symbols, authorship, currencies)
that the model demonstrably reproduces in open-ended completion
(Sec.~\ref{sec:future}); they play the role of a known-facts control:
behaviorally verified and token-matched, hence also clean, in the sense of
free of the sentence-format confounds discussed in Sec.~\ref{sec:limits}.
The external datasets are CounterFact~\cite{counterfact} and
TruthfulQA~\cite{truthfulqa} ($250$ pairs, pinned revisions in
Sec.~\ref{sec:repro}), which contain many relations a $1.5$B model plausibly
does not know.

To quantify this, we define the \emph{know-rate}: the fraction of facts the
model reproduces correctly in open-ended greedy completion. Over $200$
CounterFact prompts, Qwen2.5-1.5B has a know-rate of $30.0\%$ ($61/200$),
which grounds the ``plausibly does not know'' reading; the know-rates of the
other models appear in Sec.~\ref{sec:generalization}.


\paragraph{Why a flipped predicate and not a negation.} The false counterpart
is built by flipping a content word (``planet'' for ``star''), not by
inserting ``not''. The choice is deliberate, for three reasons. First,
negation is not a neutral edit: it activates a dedicated polarity direction in
the representation~\cite{burger2024}, so a negated counterpart would entangle
the truth signal with the polarity signal. Our design keeps both sentences
affirmative and isolates truth; polarity is then studied separately and
deliberately in the $2{\times}2$ design of Sec.~\ref{sec:poldesign}. Second,
inserting ``not'' changes the token count and shifts all subsequent positions,
weakening the token-level matching that makes the intra-pair difference
cancel shared structure. Third, the two designs are complementary rather than
competing: the negated forms are exactly the NT/NF cells of the polarity
design, so the alternative is not ignored but treated as its own object of
study.

\subsection{The axis, and where labels enter}
\label{sec:axis}

We form the intra-pair difference matrix $D$ with rows
$h_{\text{true}} - h_{\text{false}}$ at a given layer, and take its SVD,
\begin{equation}
D = U S V^{\top}, \qquad v_{1} = V_{1},
\end{equation}
where $v_1$, the top right-singular vector, is the dominant direction of the
differences. The role of supervision in this construction can be stated
exactly.

\begin{lemma}[Label-free identification of the axis, up to sign]
\label{lem:labels}
Let each pair $i$ contribute a row $d_i = s_i\,(h_i^{A} - h_i^{B})$ to $D$,
where $h_i^{A}, h_i^{B}$ are the two hidden states of the pair and
$s_i \in \{\pm 1\}$ encodes which sentence of the pair is designated true.
Then $D^{\top} D = \sum_i d_i d_i^{\top}$ does not depend on $(s_i)$; hence
the right singular vectors of $D$, and in particular the span of $v_1$, are
invariant to the within-pair labels.
\end{lemma}

\begin{proof}
$d_i d_i^{\top} = (-d_i)(-d_i)^{\top}$ for every $i$, so each summand of the
Gram matrix, and therefore the Gram matrix and its eigenvectors, is unchanged
under any sign assignment.
\end{proof}

Labels enter the pipeline at exactly two points. First, a single global
orientation bit (``the true side is positive''), fit on training folds,
fixes the sign of $v_1$. Second, the evaluation itself: AUC is defined with
respect to labels. Everything upstream of these two points, the pairing, the
differences, and the axis, uses no labels. We therefore describe the axis as
\emph{training-free} (no optimization of any kind) and \emph{label-free up to
one sign bit} (Lemma~\ref{lem:labels}), and we reserve the word
``unsupervised'' for the polarity recovery of Sec.~\ref{sec:polarity}, where
the recovered structure (polarity) is never labeled at all. The lemma also
explains, in advance, the behavior of the permutation null below: permuting
which sentence of a pair is called true leaves the span of the axis
unchanged and destroys only the orientation and the test labels, which is why
the null concentrates near chance.

\subsection{A falsifiable evaluation protocol}
\label{sec:protocol}

The central methodological point is that the axis, its calibration, the
orientation bit, \emph{and} the choice of peak layer (the layer whose held-out AUC is maximal, chosen on training folds) must be fit on training
data only. We use $k$-fold cross-validation \emph{over pairs}: a pair is
never split across folds, which would leak the shared topic. Everything is
fit on the training folds and we report held-out AUC averaged across folds.
Pairs from the same CounterFact relation template may still fall in
different folds; generalization across topics is therefore not certified by
this protocol, and we return to the point in Sec.~\ref{sec:limits}.

Part of the apparent accuracy of any layer-scanned pipeline is selection: the
best of $29$ layers looks good by chance. The \emph{permutation null}
quantifies exactly that contribution. Let $A_{\mathrm{obs}}$ be the observed
maximum held-out AUC over layers ("AUC is the area under the ROC curve, equivalently the probability that a randomly chosen true statement receives a higher score than a randomly chosen false one, with ties counted as one half."). For $b = 1, \dots, B$ we randomly swap
true/false \emph{within} each pair (the labels become noise while the pair
structure is preserved), rerun the identical pipeline including the layer
scan, and record the maximum held-out AUC $A_b$. The estimated $p$-value is
the standard conservative estimator
\begin{equation}
\hat{p} \;=\; \frac{1 + \#\{\,b : A_b \ge A_{\mathrm{obs}}\,\}}{B + 1}.
\end{equation}
With $B = 200$ and no null value reaching the observed statistic,
$\hat{p} = 1/201 \approx 0.005$, the value reported in
Table~\ref{tab:main}. As a reference point we also fit an
$\ell_2$-regularized logistic probe on the raw $d$-dimensional states, held
out over the same pair folds.

\subsection{A complex-plane extension}
\label{sec:complexmethod}

We additionally form a second axis $v_2$, the SVD residual orthogonal to
$v_1$, treat the projection as $z = c_1 + i\,c_2$ with calibrated real and
imaginary parts, and test whether magnitude- or phase-based scores improve
separation over the real axis alone. They do not (Sec.~\ref{sec:phase}); the
negative finding is kept because it is itself informative.

\subsection{Polarity design}
\label{sec:poldesign}

To test whether a genuine second dimension exists, we use a $2{\times}2$
design: each fact in four forms, affirmative-true (AT), affirmative-false
(AF), negated-true (NT), negated-false (NF), giving the XOR structure
of~\cite{burger2024}. We compute the affirmative and negated truth directions
$t_A = \overline{AT} - \overline{AF}$, $t_N = \overline{NT} - \overline{NF}$,
the general direction $t_G \propto t_A + t_N$, and the supervised polarity
direction $t_P \propto t_A - t_N$. For the recovery test, we run the
difference SVD on minimal pairs that \emph{mix} the two polarities, so the
SVD never sees the polarity label, and compare its leading directions to the
supervised $t_P$.

\subsection{Residual decomposition and its correctness gate}
\label{sec:decomp}

Within a block the residual stream is updated additively,
\begin{equation}
h_{\ell+1} \;=\; h_{\ell} + a_{\ell} + f_{\ell},
\end{equation}
where $a_{\ell}$ is the attention contribution and $f_{\ell}$ the FFN
contribution, both captured by forward hooks. The residual stream studied
everywhere else in this paper is the sum; the anatomy of
Sec.~\ref{sec:anatomy} reads the two addends separately.

Every anatomy tool embeds an \textbf{identity check}: it verifies that
$a_{\ell} + f_{\ell}$ reconstructs the residual delta $h_{\ell+1} - h_{\ell}$.
The verdict uses the \emph{median} relative error across (sentence, layer)
rather than the maximum: where a block adds almost nothing the delta is near
zero and the ratio explodes, so the maximum is fragile while the median
reports the typical case. Observed median error: about $1.0 \times 10^{-7}$
on both datasets, so the decomposition is exact and the anatomy numbers are
trustworthy. (The single inflated layer, block $27$, has a near-zero residual
delta and is excluded from interpretation.) Anatomy runs use
\texttt{float32}: the identity fails spuriously in \texttt{bfloat16} because
subtracting large residual states to form small per-block deltas suffers
catastrophic cancellation.

\section{Results I: signal, dimensionality, polarity}
\label{sec:results}

All numbers are on Qwen2.5-1.5B with last-token pooling; AUC is out-of-fold
unless stated.

\subsection{A real but bounded truth signal}
\label{sec:signal}

Table~\ref{tab:main} reports the held-out evaluation at the peak layer (layer
$16$). On the curated pairs the training-free axis reaches AUC $0.938$ with
permutation $\hat{p} = 0.005$ (null mean $0.588$, $95$th percentile $0.664$):
the geometric axis aligns with truth well beyond what layer selection alone
can manufacture. The signal is therefore real. It is also \emph{bounded}: a
full linear probe on the same raw states reaches $0.989$. The probe sees
information the single axis discards.

\begin{table}[t]
\centering
\caption{Held-out AUC ($5$-fold CV over pairs) at layer $16$, Qwen2.5-1.5B.
``1D'' is the real axis; ``MAG'' the magnitude-gated polar score; ``PHASE''
the angular-deviation score; ``2D'' a logistic on
$(\mathrm{Re},\mathrm{Im})$; ``Probe'' an $\ell_2$ logistic on raw states.
The last column is the label-permutation $\hat{p}$ for the 1D axis.}
\label{tab:main}
\begin{tabular}{lccccccc}
\toprule
Dataset & pairs & 1D & MAG & PHASE & 2D & Probe & $\hat{p}$ \\
\midrule
Curated pairs                    & $45$  & $0.938$ & $0.936$ & $0.928$ & $0.931$ & $0.989$ & $0.005$ \\
CounterFact only                 & $250$ & $0.792$ & $0.790$ & $0.771$ & $0.795$ & $0.854$ & $0.005$ \\
CounterFact / TruthfulQA mix     & $250$ & $0.638$ & $0.640$ & $0.676$ & $0.651$ & $0.841$ & $0.005$ \\
\bottomrule
\end{tabular}
\end{table}

\subsection{The gap scales with model knowledge, and with heterogeneity}
\label{sec:dim}

The rows of Table~\ref{tab:main} differ not only in difficulty but in the
\emph{size of the probe gap}, and the three rows separate two distinct
effects.

The first effect is knowledge. On curated facts the model demonstrably knows,
the axis ($0.938$) trails the full probe ($0.989$) by $0.051$. On CounterFact
alone, where the measured know-rate is $30.0\%$ (Sec.~\ref{sec:data}) and
many relations are obscure (for example, the mother tongue of a minor
historical figure), the separability of the axis falls to $0.792$ and the gap
widens to $0.062$ (probe $0.854$). Less knowledge, lower separability, and a
modestly larger gap.

The second effect is heterogeneity. On the CounterFact/TruthfulQA mix the gap
jumps to $0.203$ (axis $0.638$, probe $0.841$), far more than either
ingredient alone would suggest. A single direction fit across heterogeneous
material must serve all of it at once, while the full probe is free to
combine directions; the excess gap on the mix therefore measures how much the
truth direction itself varies with the material. Subsequent work confirms
this reading independently (Sec.~\ref{sec:generalization}, item 4): the
per-category truth axes form a structured, semantically signed geometry, and
forcing one axis onto a mixture of categories is exactly what the mix row
pays for.

We read the knowledge effect as a statement about \emph{dimensionality}. When
the model knows a fact, its truth content is concentrated and a single
direction captures most of it. When the model does not know a fact, the
truth-relevant signal is weaker and more diffuse, and one direction recovers
less of it. The dimensionality of the truth representation is thus not fixed
but \emph{knowledge-dependent}; the mix row adds the complementary statement
that it is also \emph{category-dependent}. This joint framing is, to our
knowledge, stated in these terms by neither the linear-probe nor the
dictionary-learning literature, and it is the contribution we consider most
novel.

\subsection{The complex plane does not help; ``phase'' is the sign flip}
\label{sec:phase}

On single-polarity data the second dimension adds nothing: at layer $16$ the
2D logistic ($0.931$) does not exceed the 1D axis ($0.938$) on curated pairs,
and on CounterFact it gains $+0.003$ ($0.795$ vs.\ $0.792$). This is
structural. Minimal pairs cancel the topic, so the differences concentrate
along one direction; the SVD places the entire signal in $v_1$, and $v_2$ is
residual variance, noise rather than a second truth factor. A controlled
synthetic check confirms that the apparent ``rotation'' of falsehoods is an
artifact: with a pure sign flip and zero genuine rotation, an angular (phase)
score still separates classes, because a state at $\mathrm{Re} < 0,\
\mathrm{Im} \approx 0$ has phase $\approx \pi$ by definition. The earlier
intuition of truth-as-phase has no operational content here: it reduces to
the sign of the real axis re-expressed in polar coordinates. We retract that
framing.

\subsection{Recovery of the polarity direction without polarity labels}
\label{sec:polarity}

A genuine second dimension does exist, but only when the data span both
polarities, and it is the polarity direction of~\cite{burger2024}, not a
phase. Table~\ref{tab:polarity} shows the diagnostic: an axis fit on
affirmative statements classifies affirmative held-out facts at $0.981$ but
\emph{anti-classifies} negated ones at $0.079$, a near-total flip, while the
general direction $t_G$ recovers the negated facts at $0.954$. The truth
direction genuinely rotates with polarity ($\cos(t_A, t_N) = -0.415$).

The methodological result is the recovery. When we run the difference SVD on
\emph{mixed-polarity} minimal pairs, blind to the negation label, its leading
direction aligns with the supervised polarity direction $t_P$ at cosine
$0.959$ (layer $15$), and $t_P$ lands in the first principal component, while
alignment with $t_G$ is only $0.157$. In other words, when both polarities
are present, the dominant axis of the differences \emph{is} polarity, and the
SVD finds it without supervision; by Lemma~\ref{lem:labels} the recovery
cannot have used the labels even implicitly. Where~\cite{burger2024}
constructs $t_P$ by separating affirmative from negated statements by hand,
the same direction emerges here from geometry alone.

\begin{table}[t]
\centering
\caption{Polarity analysis (leave-facts-out CV), layer $14$ chosen on
affirmative data; and unsupervised $t_P$ recovery at layer $15$. AUC below
$0.5$ indicates a sign flip.}
\label{tab:polarity}
\begin{tabular}{lc}
\toprule
Quantity & Value \\
\midrule
$\cos(t_A, t_N)$ & $-0.415$ \\
Affirmative axis on affirmative (held-out AUC) & $0.981$ \\
Affirmative axis on \emph{negated} (held-out AUC) & $0.079$ \\
General direction $t_G$ on negated (held-out AUC) & $0.954$ \\
\midrule
SVD vs.\ supervised $t_P$: $\max_k |\cos(\mathrm{SVD}_k, t_P)|$ & $0.959$ \\
\quad (principal component carrying $t_P$) & PC\,$1$ \\
$\cos(\mathrm{SVD}_1, t_G)$ & $0.157$ \\
\bottomrule
\end{tabular}
\end{table}


\section{Robustness: seven falsification experiments}
\label{sec:robust}

The first half of this paper established a positive claim. We now stress-test
it. The claim under test is:

\begin{quote}
\emph{In Qwen2.5-1.5B, the truth value of a statement is carried by a single
(one-dimensional) direction, the position axis obtained from the SVD of
paired true/false differences, on facts the model knows; the second and
higher dimensions add no separation.}
\end{quote}

A claim of one-dimensionality is falsifiable in a specific way: exhibit a
second, third, or depth-distributed coordinate that separates true from false
better than the single axis. We construct seven such coordinates and test
each.

\paragraph{Two kinds of test, kept distinct.} We separate \emph{controlled
experiments} from \emph{visual observations}, and label every result as one
or the other. Controlled experiments (experiments 1--4) produce a held-out
AUC under $k$-fold cross-validation over pairs and a label-permutation null
that quantifies selection bias; a result is meaningful only if it clears that
null. Visual observations (experiments 5--7) do not classify; they project
the states and measure a single descriptive statistic, the per-coordinate
separation AUC or the centroid-separation ratio, so the conclusion follows
from a number attached to a picture, not from a fitted model. Throughout, the
model is frozen Qwen2.5-1.5B, last-token pooling, peak layer $16$ unless
stated.

Table~\ref{tab:seven} summarizes all seven with their real numbers; the prose
below gives the reasoning for each.

\begin{table}[t]
\centering
\caption{Seven experiments attempting to falsify the one-dimensional truth
axis. ``Exp.'' = controlled experiment (held-out CV + permutation null);
``Obs.'' = visual observation (descriptive separation statistic). None
overturns the claim.}
\label{tab:seven}
\small
\begin{tabular}{p{0.5cm}p{4.6cm}p{1.1cm}p{4.6cm}p{1.4cm}}
\toprule
\# & Richer structure hypothesized & Type & Result (real numbers) & Verdict \\
\midrule
1 & A complex/phase second dimension (truth as rotation) & Exp. & PHASE AUC $0.928$ vs.\ 1D $0.938$ (curated); reduces to the real-axis sign flip & falsified \\
2 & Truth is global across depth (fixed-axis signature) & Exp. & global $0.923$ vs.\ local $0.938$ ($-0.015$), $\hat{p} = 0.0099$; persistence gap $+0.04$ & falsified \\
3 & Truth is dynamic, a ``butterfly effect'' in the unsaturated shade, on hard data & Exp. & global $0.642$ vs.\ local $0.627$ ($+0.015$), $\hat{p} = 0.0099$; commitment true $2.83 \to 5.91$ vs.\ false $2.65 \to 5.86$ & falsified \\
4 & Truth is distributed across layer ``departments'' & Exp. & diagnostic curve is a smooth ramp (no steps); global $0.646$ vs.\ local $0.627$ ($+0.019$), $\hat{p} = 0.0099$ & falsified \\
5 & The argument $\theta$ (phase) carries truth & Obs. & $\theta$ AUC $0.504$ / $0.584$ / $0.556$ (CF / TQA / mix) & falsified \\
6 & The energy $|z|$ (commitment) carries truth & Obs. & $|z|$ AUC $0.560$ / $0.558$ / $0.543$ & falsified \\
7 & True/false form distinct 3D clumps & Obs. & centroid ratio $0.55$ / $0.51$ / $0.48$ ($<1$, overlapping); $\mathrm{Re}$ alone $0.795$ / $0.656$ & falsified \\
\bottomrule
\end{tabular}
\end{table}

\paragraph{1. The complex plane.} The original framing promoted the second
SVD direction to an imaginary axis and read truth as a phase rotation. If
real, a phase- or magnitude-based score should beat the real axis. It does
not: on curated pairs the phase score reaches AUC $0.928$ against $0.938$ for
the real axis, and a controlled synthetic check shows that with a pure sign
flip and \emph{zero} genuine rotation an angular score still separates,
because a state at $\mathrm{Re} < 0$ has phase $\approx \pi$ by definition.
The ``rotation'' is the sign of the real axis re-expressed in polar
coordinates. The complex framing is decorative; we retract it. (This is the
negative result of Sec.~\ref{sec:phase}, restated as the first falsification
experiment.)

\textbf{2. A global depth signature.} If truth flows through the network, the full trajectory $[\Delta_0, \dots, \Delta_{28}]$ should beat the single best layer, where the signed deviation at each layer $\ell$ is defined as $\Delta_\ell = b_\ell - \frac{1}{2}$, with $b_\ell = \sigma(\operatorname{Re}(z_\ell))$ representing the calibrated projection passed through a sigmoid along a fixed axis. On curated pairs it does not: the depth-signature classifier reaches held-out AUC 0.923 against 0.938 for the single layer (-0.015), and although the signature is real (permutation $\hat{p} = 0.0099$), it carries no advantage. The crossing-persistence gap between false and true is $+0.04$, flat (where for each sentence, the first layer at which the trajectory crosses below zero and the fraction of subsequent layers it stays below; the gap is the difference of this persistence between false and true statements). No propagation.


\paragraph{3. A dynamic butterfly effect in the shade.} The signed deviation
saturates ($\sigma$ compresses large values to $\pm\tfrac12$), so a
propagating \emph{uncertainty} could be invisible to it. We re-ran on hard
data (CounterFact/TruthfulQA mix) with the unsaturated features
$[\mathrm{Re}, |z|]$, the ``shade''. The depth signature now beats the single
layer, but by $+0.015$ ($0.642$ vs.\ $0.627$, $\hat{p} = 0.0099$): real and
marginal. The decisive descriptive measure kills the butterfly effect: the
commitment magnitude grows identically for both classes (true
$2.83 \to 5.91$, false $2.65 \to 5.86$); false statements do not hesitate
more than true ones, and there is no propagating divergence.

\paragraph{4. Layer departments.} If different layer bands perform different
jobs (surface, semantics, logic), a single semantic axis re-read everywhere
wastes most layers; giving each layer its \emph{own} native probe should
recover a distributed signal, and the per-layer separation curve should show
steps. It shows a smooth ramp instead: the per-layer native separation rises
continuously to a single peak at layer $16$, with no plateaus and no jumps.
The per-layer signature beats the single layer by $+0.019$ ($0.646$ vs.\
$0.627$, $\hat{p} = 0.0099$), again real and marginal. No departments; the
native-axis ``cure'' improves the fixed axis but does not change the ceiling.

\paragraph{5--6. Energy and argument.} Decomposing each state into position
($\mathrm{Re}$), energy ($|z|$), and argument ($\theta$) and measuring how
well each separates true from false (Fig.~\ref{fig:coords}): only position
carries signal. The energy separates at $|z|$ AUC $0.560$/$0.558$/$0.543$
(CounterFact/TruthfulQA/mix): true and false are equally committed. The
argument separates at $\theta$ AUC $0.504$/$0.584$/$0.556$, essentially
chance; what little it shows is the echo of the $\mathrm{Re}$ sign, not
independent information. The phase, the very coordinate the complex framing
made central, is empty.

\paragraph{7. Three-dimensional clumping.} If truth were a richer object,
true and false should form distinct clusters in the three-dimensional
$(v_1, v_2, v_3)$ space. They do not (Fig.~\ref{fig:plane}). The
centroid-separation ratio, the distance between the true and false centroids
divided by the average cloud spread, is $0.55$/$0.51$/$0.48$, all below $1$:
the clouds overlap. Crucially this low ratio is \emph{not} evidence against
the axis; it is the one-dimensional signal seen through two noise dimensions.
Projected onto position alone, the same data separate at AUC $0.795$
(CounterFact). Adding the orthogonal dimensions dilutes, it does not enrich.

\subsection{What survives}
\label{sec:survive}

Three statements resist all seven experiments.

\emph{Truth is one-dimensional and lives on the position axis.} Across every
decomposition, only $\mathrm{Re}$ separates; energy, phase, the second and
third SVD directions, and the full-depth trajectory add at most $+0.02$ AUC
and usually nothing.

\emph{Separability scales with knowledge.} On CounterFact, where the model
often knows the fact, the position axis separates at AUC $0.795$; on the more
ambiguous TruthfulQA at $0.656$; on the diffuse mix the centroids nearly
coincide. The clearer the model's knowledge, the more one-dimensional and
separable the truth representation.

\emph{On unknown facts, the classes show low separability.} Where the model
does not know the fact, true and false representations do not form distinct
clusters but overlap in space, intuitively a dense, uniform distribution
rather than two distinct clumps. We state this as \emph{low class
separability / distribution overlap} (centroid ratio $<1$), the measured
counterpart of the knowledge-dependent dimensionality of Sec.~\ref{sec:dim}.

\begin{figure}[t]
\centering
% Upload counterfact_coords.png and truthfulqa_coords.png alongside the .tex.
\begin{minipage}{0.49\textwidth}\centering
\includegraphics[width=\textwidth]{counterfact_coords.png}\\[2pt]
{\footnotesize (a) CounterFact: $\mathrm{Re}$ $0.795$, $|z|$ $0.560$, $\theta$ $0.504$}
\end{minipage}\hfill
\begin{minipage}{0.49\textwidth}\centering
\includegraphics[width=\textwidth]{truthfulqa_coords.png}\\[2pt]
{\footnotesize (b) TruthfulQA: $\mathrm{Re}$ $0.656$, $|z|$ $0.558$, $\theta$ $0.584$}
\end{minipage}
\caption{Experiments 5--6 and knowledge scaling, layer $16$. Per-coordinate
separation AUC. On both datasets only the position axis ($\mathrm{Re}$)
separates true from false; energy ($|z|$) and argument ($\theta$) are at
chance. The position separation falls from $0.795$ (CounterFact, often-known
facts) to $0.656$ (TruthfulQA, more ambiguous), while the one-dimensional
structure is unchanged.}
\label{fig:coords}
\end{figure}

\begin{figure}[t]
\centering
% Upload counterfact_plane.png and counterfact_3d.png alongside the .tex.
\begin{minipage}{0.49\textwidth}\centering
\includegraphics[width=\textwidth]{counterfact_plane.png}
\end{minipage}\hfill
\begin{minipage}{0.49\textwidth}\centering
\includegraphics[width=\textwidth]{counterfact_3d.png}
\end{minipage}
\caption{Experiment 7, CounterFact at layer $16$. Left: the 2D plane; the
true (blue) and false (red) centroids (diamonds) are close relative to the
cloud spread (centroid-separation ratio $0.55$, overlapping). Right: the
three-dimensional $(v_1, v_2, v_3)$ scatter shows a gradient along $v_1$
(position) but no distinct clumps. The separation is confined to the
horizontal position axis; the extra dimensions are noise.}
\label{fig:plane}
\end{figure}

\begin{figure}[t]
\centering
% Upload truthfulqa_plane.png and truthfulqa_3d.png alongside the .tex.
\begin{minipage}{0.49\textwidth}\centering
\includegraphics[width=\textwidth]{truthfulqa_plane.png}
\end{minipage}\hfill
\begin{minipage}{0.49\textwidth}\centering
\includegraphics[width=\textwidth]{truthfulqa_3d.png}
\end{minipage}
\caption{Experiment 7, TruthfulQA at layer $16$. The same picture on more
ambiguous facts: centroid-separation ratio $0.51$ (overlapping), no
three-dimensional clumps. Lower knowledge, lower separability, identical
one-dimensional structure.}
\label{fig:plane_tqa}
\end{figure}


\section{Results II: the anatomy of the axis}
\label{sec:anatomy}

The falsification battery established \emph{that} truth is one-dimensional on
known facts. This section asks \emph{how the residual stream comes to carry
it}: which architectural component, attention or the FFN, produces, holds, or
degrades the signal, layer by layer. The results here are Tier 2 in the sense
of Sec.~\ref{sec:scope}: same protocol and shields, but smaller margins in
places, and interpretations are marked as such. The decomposition and its
correctness gate are defined in Sec.~\ref{sec:decomp}.

\subsection{Per-layer readability: a stable alternation}
\label{sec:alternation}

For each layer we fit the truth axis on $a_{\ell}$ and on $f_{\ell}$
separately and measure held-out AUC, repeating over $5$ cross-validation
seeds. A per-layer difference $\mathrm{AUC}(a_\ell) - \mathrm{AUC}(f_\ell)$
is counted as real only if its magnitude exceeds both $2\sigma$ (its own
across-seed standard deviation) and a floor of $0.03$. Table~\ref{tab:peak}
shows the peak block.

\begin{table}[h]
\centering
\caption{Truth AUC by component at the readability peak (block $15$),
CounterFact, mean $\pm$ std over $5$ seeds. Residual = the summed stream
(reference).}
\label{tab:peak}
\begin{tabular}{lccc}
\toprule
 & residual & attention & FFN \\
\midrule
AUC & $0.792 \pm 0.003$ & $0.845 \pm 0.003$ & $0.734 \pm 0.004$ \\
\bottomrule
\end{tabular}
\end{table}

The alternation survives re-seeding and reproduces structurally across both
datasets. On CounterFact, attention stably dominates readability at layers
$\{6,7,10,12,13,14,15,17,18\}$ and the FFN at
$\{8,11,16,19,20,21,22,23,25,26\}$. At the peak (block $15$), where the truth
axis is read, attention leads by $+0.110 \pm 0.004$. \textbf{Interpretation
(not yet causal):} attention appears to carry the truth signal most strongly
through the middle band, including the peak. Readability, however, is not
necessity, which Sec.~\ref{sec:ablation} tests directly, and which changes
the picture.

\subsection{The expanded FFN space does not help}
\label{sec:expanded}

The FFN expands its input to $8960$ dimensions, applies SwiGLU, and
re-compresses it. One might expect the truth signal to be more cleanly
readable inside that larger, putatively sparser space, before compression. We
capture the expanded activation (the input to \texttt{down\_proj}) and fit
the truth axis there. In this $p \gg n$ regime the permutation null is
elevated, so only the margin over the null counts: at the peak the expanded
axis clears its null cleanly (curated $0.953$ vs.\ null $95$th percentile
$0.669$, margin $+0.284$, $\hat{p} = 0.0099$; CounterFact $0.816$ vs.\
$0.561$, margin $+0.254$, $\hat{p} = 0.0099$). \textbf{But relative to the
compressed residual at the same layer, the expanded space gains essentially
nothing}: $+0.017$ (curated) and $+0.023$ (CounterFact). Raw architectural
expansion does not undo superposition on these facts. The implication for
future work is concrete: sparsity must be \emph{imposed}, for example by a
trained sparse autoencoder, not merely afforded by a wider space.

\subsection{Causal ablation: attention builds, the FFN erodes}
\label{sec:ablation}

To separate readability from necessity we zero a component's contribution
over a band of layers (a hook returns zeros, so $h + 0 = h$), let the effect
propagate, and measure the collapse of truth AUC at a downstream readout
(Table~\ref{tab:abl}).

\begin{table}[h]
\centering
\caption{Causal ablation on CounterFact ($250$ pairs). Drop $=$ intact $-$
ablated; positive drop $=$ the component is causally necessary.}
\label{tab:abl}
\begin{tabular}{llccc}
\toprule
band (layers) & readout & intact & attention-off (drop) & FFN-off (drop) \\
\midrule
$11$--$15$ (middle)     & block $15$ & $0.793$ & $0.487\;(+0.306)$ & $0.819\;(-0.026)$ \\
$16$--$18$ (transition) & block $18$ & $0.712$ & $0.708\;(+0.004)$ & $0.815\;(-0.103)$ \\
$19$--$26$ (late)       & block $27$ & $0.537$ & $0.534\;(+0.003)$ & $0.539\;(-0.002)$ \\
\bottomrule
\end{tabular}
\end{table}

Three readings, in order of confidence:
\begin{itemize}
\item \textbf{Middle band: attention is causally necessary.} Removing
attention collapses truth from $0.793$ to $0.487$ (drop $+0.306$); removing
the FFN does nothing ($-0.026$). Attention does not merely read the signal
here, it builds it.
\item \textbf{Transition band: the FFN \emph{degrades} the signal.} Removing
attention again does nothing ($+0.004$), but removing the FFN \emph{raises}
truth AUC from $0.712$ to $0.815$ (drop $-0.103$). From layer $16$ onward the
FFN injects content that lowers true/false separability. This is the first
causal support for an earlier conjecture, that the FFN adds context which
amplifies the overlap on uncertain facts, though it rests on one band and one
dataset and should be treated as a lead, not a result.
\item \textbf{Late band: nothing remains to destroy.} The intact AUC is
already about $0.54$ (near chance): the truth signal has evaporated toward
the output, so neither ablation has an effect.
\end{itemize}

\textbf{This falsifies the initial ``attention builds, FFN consolidates''
hypothesis.} There is no late consolidation; the FFN's apparent late
\emph{readability} dominance (Sec.~\ref{sec:alternation}) was the readability
of a dying signal, not active maintenance. The corrected picture: attention
constructs the truth representation in the middle layers; the FFN begins to
erode it in the transition band; by the late layers it is gone, as the model
shifts from representing the fact to predicting the next token.

\section{Discussion: what a single axis can and cannot read}
\label{sec:discussion}

It is tempting to view a cheap geometric probe as a low-cost version of
dictionary learning, reading superposition without expanding it. We argue
this is impossible, and that the impossibility is a capacity statement, not
an implementation gap. A projection is an inner product: it collapses a
$d$-dimensional state to a single scalar, discarding $d - 1$ dimensions. From
that scalar one cannot reconstruct which features were present, only how much
a chosen direction was. ``Reading the superposition'' requires asking the
question once per dictionary atom, which over an overcomplete dictionary of
size $m$ is exactly the $O(d\,m)$ sparse autoencoder. Compression (project to
a scalar) and disambiguation (expand to a sparse code) are not inverse
cheap-vs-expensive versions of one operation; they are different operations,
and only the second reads superposition. The gaps of
Sec.~\ref{sec:dim}, $0.062$ on homogeneous unknown material and $0.203$ on
the heterogeneous mix, are direct measurements of the truth information that
a single axis structurally cannot return.

The consequence is a clear scope. This method's value is \emph{structural and
cheap}: a training-free axis that, on facts the model knows, rivals a trained
probe at two dot products per token; a quantitative law relating the
readability of truth to the model's knowledge of the fact; and a route to the
polarity direction that needs no polarity labels. It is not a substitute for
sparse autoencoders and should not be presented as one.

\section{The axis reads truth and confidence}
\label{sec:meta}

The seven experiments establish that there is no second coordinate of truth
to find. This has a constructive consequence that is easy to miss, so we
state it plainly. The same axis admits \emph{two} readings, not one.

Its \emph{sign} classifies a statement: $\mathrm{Re} > 0$ is the true side,
$\mathrm{Re} < 0$ the false side. Its \emph{separability}, how cleanly true
and false split on a given set of facts, measures how well the model knows
that material. The geometry that answers ``is this statement true?''
therefore also answers, indirectly, ``does the model know this?'', because a
model that does not know produces overlapping rather than separated
representations (Figs.~\ref{fig:plane},~\ref{fig:plane_tqa}). These are not
two axes; they are two readings of one axis: the direction of the arrow, and
the sharpness of the arrow.

The reading is quantitative at the level of a distribution. The position-axis
separation is $0.795$ on CounterFact and $0.656$ on TruthfulQA: by this
measure the model ``knows'' the CounterFact material better than the
TruthfulQA material, read off the geometry alone, with no external grading.
We report this as a property of the dataset, not yet of the single fact.
Calibrating a per-fact confidence, a local separability around one statement
that predicts whether the model actually errs on it, is the natural next
step, and requires a behavioral label (does the model answer the open
question correctly?) that the minimal-pair design does not contain. We mark
it as future work rather than a claim. Section~\ref{sec:future} returns to
this label and to a prerequisite it exposed.

\section{Why unknown facts overlap: the signature of superposition}
\label{sec:why}

It remains to interpret \emph{why} the classes overlap on unknown facts
rather than simply asking \emph{that} they do. The overlap is not measurement
noise; it is the observable signature of unresolved superposition, and the
explanation follows from the capacity argument of Sec.~\ref{sec:discussion}.

A projection onto a single direction returns a scalar: it collapses the
$d$-dimensional state to one number and discards $d - 1$ dimensions. When the
model knows a fact, its truth content has been allocated a clean, dominant
direction of its own (the residual stream has, in effect, dedicated an axis
to it), and the single projection recovers almost all of it (position AUC
$0.795$, near the full-probe ceiling). When the model does not know a fact,
there is no dedicated direction: the weak truth-relevant signal remains
distributed across many superposed features, and no single axis can pull it
out of the tangle. Recovering it would require asking the question once per
dictionary atom, the expansion that a sparse autoencoder performs at cost
$O(d\,m)$, which is a different operation from a single projection, not a
cheaper version of it.

The consequence is exactly what the visual observations show. On known facts
the truth occupies its own direction, the classes form two regions, and the
centroid ratio is high. On unknown facts the truth is spread across the
superposition, the classes overlap, and the centroid ratio falls below one.
The dense, uniform overlap of true and false on unknown facts is therefore
not a failure of the probe but a direct, visual readout of superposition that
the probe structurally cannot resolve. The probe gaps of
Sec.~\ref{sec:dim} and the low centroid ratio measured here are views of the
same quantity: the truth information that lives off-axis on facts the model
does not know ($0.062$), inflated further when heterogeneous material forces
one axis to serve many truth frames at once ($0.203$). Subsequent work
(Sec.~\ref{sec:generalization}, item 4) begins to resolve that superposed
remainder into named, signed components: what this paper can only call
unresolved superposition acquires, there, a geometry.

\section{Scope and generalization: what subsequent work established}
\label{sec:generalization}

Everything above concerns Qwen2.5-1.5B. The natural question, whether the
same conclusions hold on other models, has been addressed in subsequent work
by the author, available in the same repository, and we state the outcomes
here as precise claims so the reader knows what is established and where.

\begin{enumerate}
\item \textbf{The signal and the peak structure replicate on four models
across two families} (Qwen2.5-1.5B, Qwen2.5-3B, Llama-3.2-1B, Llama-3.2-3B):
on each model the training-free axis clears its permutation null held-out,
with a single readability peak at intermediate relative depth.
\item \textbf{The knowledge-gap relation gains cross-family support, with a
family offset.} With knowledge measured behaviorally and uniformly (the
know-rate of Sec.~\ref{sec:data}: $200$ CounterFact prompts, greedy decoding,
seed $0$) and all gaps measured on CounterFact under the same held-out
protocol, the four models give Table~\ref{tab:knowgap}. Within each family
the gap decreases as the know-rate increases (Qwen: $0.062$ at $30.0\%$ vs.\
$0.049$ at $37.0\%$; Llama: $0.141$ at $35.5\%$ vs.\ $0.112$ at $41.0\%$).
Across families a family-level offset remains: at comparable knowledge, the
Qwen models concentrate more of the truth signal on the single axis than the
Llama models. The knowledge dependence of Sec.~\ref{sec:dim} is therefore not
an artifact of one model, and the family offset is consistent with the
family-trait finding of item 3.
\begin{table}[H]
\centering
\caption{Behavioral know-rate ($200$ CounterFact prompts, greedy decoding,
seed $0$) and axis-vs-probe gap (CounterFact, $250$ pairs, held-out), per
model. The Qwen2.5-1.5B row is measured directly in this paper
(Table~\ref{tab:main}); the other rows come from the companion work under the
same protocol.}
\label{tab:knowgap}
\begin{tabular}{lcc}
\toprule
Model & Know-rate (correct/total) & Axis-vs-probe gap \\
\midrule
Llama-3.2-3B & $41.0\%$ ($82/200$) & $0.112$ \\
Qwen2.5-3B   & $37.0\%$ ($74/200$) & $0.049$ \\
Llama-3.2-1B & $35.5\%$ ($71/200$) & $0.141$ \\
Qwen2.5-1.5B & $30.0\%$ ($61/200$) & $0.062$ \\
\bottomrule
\end{tabular}
\end{table}
\item \textbf{The anatomy generalizes as a relational law.} An intermediate
result (that the FFN writes pro-truth into the peak layer and flips to
anti-truth at exactly the layer after the peak) was published by the author
and then corrected within days by an axis-provenance control built against
the author's own instrument; the correction is reported as a result. The
corrected law, confirmed by pre-registered predictions on all four models:
the FFN's write at a block opposes every truth frame that does not contain
it, while attention aligns with existing frames it did not write. At the
peak, on the same block and the same frame, the two contributions are near
mirror images ($+1.63$ for attention vs.\ $-1.43$ for the FFN on
Llama-3.2-3B). The erosion reported in Sec.~\ref{sec:ablation} survives the
correction and is causally attributed, by exact algebraic split and by
freezing interventions, to the SwiGLU value stream rather than the gate.
\item \textbf{The axis is a semantically signed mixture.} Fitting the axis
per fact category shows that per-category truth axes form a structured
geometry (related categories cluster with positive cosines, one category
anti-aligns) and that this arrangement converges across the two model
families (Mantel test, a permutation test for the correlation between two distance
matrices, applied to the category-geometry matrices of the two families:
$r = +0.77$, $p = 0.0009$). Within-category held-out AUC exceeds
cross-category AUC, which bears directly on the topic-generalization question
raised by the protocol of Sec.~\ref{sec:protocol}.
\end{enumerate}

These four items carry their own held-out protocols, permutation nulls, and
pre-registered predictions, and are documented, with reproduction commands,
in the same repository as this paper. They are summarized here for scope, not
claimed as results of this paper. Read in sequence, the chain is linear: the
present paper isolates the axis and measures its limits; the follow-up work
explains those limits, first mechanically (which component builds the axis
and which erodes it, generalized as a relational law) and then geometrically
(what the axis is a mixture of). The overlap that Sec.~\ref{sec:why} can only
call unresolved superposition is, by the end of that chain, partially
resolved into named, signed components.

\section{A note for future work: verify behavior before geometry}
\label{sec:future}

Two natural extensions, a per-fact behavioral confidence label
(Sec.~\ref{sec:meta}) and a study of the geometry of \emph{logical} (as
opposed to semantic) truth, share a prerequisite that is easy to overlook,
and we record it here as methodological guidance rather than as a result.

\paragraph{An informal behavioral check (not a controlled experiment).}
Before building any geometric study of a capability, one should confirm that
the capability exists behaviorally. We ran a small, qualitative probe of the
base model in open-ended completion (greedy decoding, no context carried
between prompts). Two things were apparent on a handful of prompts. First,
the model reproduces the common semantic facts used throughout this paper
(capitals, chemical symbols, authorship, currencies), consistent with the
``facts the model knows'' assumption behind the curated pairs. Second, on
symbolic transitive syllogisms built from nonce words (``a \emph{grok} is
heavier than a \emph{blarn}; a \emph{blarn} is heavier than a \emph{fendu};
therefore $\ldots$''), the base $1.5$B model does \emph{not} perform the
deduction: it deflects into multiple-choice or code scaffolding rather than
completing the transitive conclusion. We stress that this is an informal
observation on a few prompts, not a measured result, and we draw no
quantitative claim from it.

\paragraph{The principle.} Studying the geometry of a capability presupposes
the capability. Measuring the hidden-state geometry of a task the model
cannot perform yields noise, and that noise is easily mistaken for ``the
property is distributed/complex'' when the truthful reading is ``the property
is absent.'' A cheap behavioral feasibility check should therefore precede
any geometric study of \emph{logical} truth; on this model that study is not
feasible, which is itself useful to know in advance. The same principle
applies to transfer: before studying the geometry of truth in a second model,
one should first verify behaviorally that the second model knows the facts in
question. (In the subsequent multi-model work of
Sec.~\ref{sec:generalization} this principle is applied systematically: the
behavioral know-rate is measured before any geometry.)

\paragraph{A hypothesis to test, not a claim.} The way the base model
deflects on uncertain input, falling back on memorized \emph{format}
scaffolding (exam layouts, question/answer templates, code), suggests, but
does not establish, that part of the overlap on unknown facts may carry a
sentence-format component rather than pure superposition noise. This is
consistent with the format confound that makes mixing datasets inadvisable
(Sec.~\ref{sec:limits}). We flag it as a direction for a future, properly
shielded experiment, not as evidence.

\section{Reproducibility}
\label{sec:repro}

Every number in this paper is reproduced by one command. \texttt{truth\_probe.py}
runs the controlled experiments (the positive results and falsification
experiments 1--4); \texttt{canvas.py} renders the visual observations
(experiments 5--7); \texttt{anatomy.py}, \texttt{anatomy\_consolidate.py},
\texttt{anatomy\_expanded.py}, and \texttt{ablation.py} produce the anatomy of
Sec.~\ref{sec:anatomy} and embed the identity, null, and stability shields
described there. Each result in \texttt{truth\_probe.py} is an \emph{isolated}
subcommand: running one task loads, extracts, and computes only its own data,
so a reviewer who wants to verify one figure runs exactly one command. All
scripts load the model with safetensors only and load external datasets as
Parquet only, never with remote code. A CUDA GPU is optional.

All scripts are archived on Zenodo (DOI 10.5281/zenodo.20938285) \\ and  mirrored
at \texttt{github.com/Francesco-Marhel/TruthProbe}.

\paragraph{Setup.}
\begin{verbatim}
pip install torch transformers datasets
\end{verbatim}

\paragraph{Positive results (Tables~\ref{tab:main},~\ref{tab:polarity}).}
\begin{verbatim}
# Table 1, top row (curated pairs): held-out AUC, ablation,
#                   full-probe baseline, permutation p-value
python truth_probe.py signal --dataset builtin --with-2d --baseline --perm 200

# Table 1, bottom row (CounterFact / TruthfulQA mix)
python truth_probe.py signal --dataset mix --max-pairs 250 \
       --with-2d --baseline --perm 200

# Table 2, top (affirmative -> negated flip; the 2D polarity structure)
python truth_probe.py polarity

# Table 2, bottom (SVD recovers the polarity direction t_P)
python truth_probe.py recovery
\end{verbatim}

\paragraph{Know-rate (Secs.~\ref{sec:data} and~\ref{sec:generalization}) and the CounterFact row of Table~\ref{tab:main}.}
\begin{verbatim}
python behav_check.py --model Qwen/Qwen2.5-1.5B --n 200 --seed 0

python truth_probe.py signal --dataset counterfact --max-pairs 250 --baseline \
       --perm 200 --rev-counterfact c945b082ca08d0a8f3ba227fb78404a09614c36e
\end{verbatim}
The same two commands with \texttt{--model} set to \texttt{Qwen/Qwen2.5-3B},
\texttt{meta-llama/Llama-3.2-1B}, or \texttt{meta-llama/Llama-3.2-3B}
reproduce the corresponding rows of Table~\ref{tab:knowgap}.

\paragraph{Controlled falsification experiments (Table~\ref{tab:seven}, 1--4).}
\begin{verbatim}
# experiment 1 (PHASE ablation) is the PHASE column of `signal` above

# experiment 2: global depth signature (fixed axis), curated pairs
python truth_probe.py domino --dataset builtin --perm 100

# experiment 3: dynamic shade on hard data
python truth_probe.py domino --dataset mix --max-pairs 250 --signal-mag --perm 100 \
       --rev-counterfact c945b082ca08d0a8f3ba227fb78404a09614c36e \
       --rev-truthfulqa 741b8276f2d1982aa3d5b832d3ee81ed3b896490

# experiment 4: per-layer native axes + diagnostic ramp
python truth_probe.py domino --dataset mix --max-pairs 250 --per-layer-axis --perm 100 \
       --rev-counterfact c945b082ca08d0a8f3ba227fb78404a09614c36e \
       --rev-truthfulqa 741b8276f2d1982aa3d5b832d3ee81ed3b896490
\end{verbatim}

\paragraph{Visual observations (Figs.~\ref{fig:coords}--\ref{fig:plane_tqa}, experiments 5--7).}
\begin{verbatim}
python canvas.py --dataset counterfact --layer 16 --out counterfact \
       --rev-counterfact c945b082ca08d0a8f3ba227fb78404a09614c36e
python canvas.py --dataset truthfulqa --layer 16 --out truthfulqa \
       --rev-truthfulqa 741b8276f2d1982aa3d5b832d3ee81ed3b896490
\end{verbatim}
Each \texttt{canvas.py} run writes \texttt{<out>\_coords.png},
\texttt{<out>\_plane.png}, and \texttt{<out>\_3d.png}. Run CounterFact and
TruthfulQA \emph{separately}; do not pass \texttt{--dataset mix} to
\texttt{canvas.py}, as mixing injects a spurious sentence-format axis.

\paragraph{Anatomy (Tables~\ref{tab:peak}--\ref{tab:abl}).}
\begin{verbatim}
# per-layer attention/FFN readability (identity check embedded)
python anatomy.py --dataset counterfact --max-pairs 250 --perm 100 \
       --rev-counterfact c945b082ca08d0a8f3ba227fb78404a09614c36e

# stability of the alternation across 5 seeds
python anatomy_consolidate.py --dataset counterfact --max-pairs 250 --seeds 5 \
       --rev-counterfact c945b082ca08d0a8f3ba227fb78404a09614c36e

# the truth axis inside the expanded FFN space (8960-dim)
python anatomy_expanded.py --dataset counterfact --max-pairs 250 --perm 100 \
       --rev-counterfact c945b082ca08d0a8f3ba227fb78404a09614c36e

# causal ablation, the three bands of Table 5
python ablation.py --dataset counterfact --max-pairs 250 \
       --band-start 11 --band-end 15 --readout 15 \
       --rev-counterfact c945b082ca08d0a8f3ba227fb78404a09614c36e
python ablation.py --dataset counterfact --max-pairs 250 \
       --band-start 16 --band-end 18 --readout 18 \
       --rev-counterfact c945b082ca08d0a8f3ba227fb78404a09614c36e
python ablation.py --dataset counterfact --max-pairs 250 \
       --band-start 19 --band-end 26 --readout 27 \
       --rev-counterfact c945b082ca08d0a8f3ba227fb78404a09614c36e
\end{verbatim}
Anatomy runs use \texttt{float32} (Sec.~\ref{sec:decomp}); about $6$ GB of
VRAM is enough for Qwen2.5-1.5B.

\paragraph{Data provenance and safety.} The curated pairs and the
$2{\times}2$ polarity set are embedded in \texttt{truth\_probe.py}. The
external datasets are \texttt{NeelNanda/counterfact-tracing} (commit
\texttt{c945b08\ldots}) and \texttt{truthful\_qa}, \emph{generation} config
(commit \texttt{741b827\ldots}), loaded as Parquet only. For bit-exact
reproduction the dataset commits may be pinned as shown, or a local file
passed with \texttt{--file-counterfact <path>}. Add \texttt{-h} to any
subcommand for its options.

\section{Limitations}
\label{sec:limits}

These results are confined to a single, small model (Qwen2.5-1.5B); the
entire battery, the positive results and all seven falsification experiments,
should be re-run at larger scale before the one-dimensionality is treated as
general (the partial replications available so far are summarized, with their
scope, in Sec.~\ref{sec:generalization}). It is plausible that departments,
dynamics, or a genuine second dimension emerge with scale, precisely because
a $1.5$B residual stream may be too compact to separate them. The datasets
are small (tens to a few hundred pairs), so confidence intervals are wide
despite the permutation control. Cross-validation is over pairs, not over
topics: pairs from the same CounterFact relation may fall in different folds,
so generalization to entirely unseen topics is not certified by the protocol
of Sec.~\ref{sec:protocol}; the within- vs.\ cross-category transfer
measurements of subsequent work (Sec.~\ref{sec:generalization}, item 4)
address the question directly. The ``unknown-fact'' interpretation of the
CounterFact gap and of the low separability is plausible but confounded: a
fact the model does not know and a fact whose truth is genuinely
multidimensional are not separated here, and the low-separability reading is
further confounded with sentence-format effects, which is why mixing datasets
must be avoided. Disentangling ignorance from genuine multidimensionality
(for example by filtering CounterFact to facts the model can answer, using
the behavioral label of Sec.~\ref{sec:future}) is necessary to make the
dimensionality claim quantitative. The anatomy results are Tier 2: margins in
the expanded-space and transition-band analyses are modest and based on
single datasets, the factual-knowledge label is noisy, and ablation measures
necessity for \emph{linear readability} of truth, not for the model's
behavior. We use last-token pooling; other aggregations may shift the axis.

\section{Conclusion}

Stripped of a decorative complex-plane framing, the positive results survive
on a real model under a falsifiable protocol: a training-free truth axis,
identified without labels up to one global sign, carries a permutation-tested
signal at cost $O(d)$; the gap between that axis and a full probe scales with
how well the model knows the fact, suggesting a knowledge-dependent
dimensionality of truth; and the difference SVD recovers the supervised
polarity direction of~\cite{burger2024} at cosine $0.959$ without polarity
labels. The one-dimensional reading withstands seven falsification
experiments, spanning a second dimension, depth-global structure, dynamics,
departments, energy, phase, and three-dimensional clustering, none of which
exhibits a richer separating coordinate. Beneath the axis, an exact
decomposition of the residual stream locates its origin: attention builds the
signal in the middle layers, causally, while the FFN erodes it after the
peak, and by the late layers the model has moved on from representing the
fact to predicting the next token. We read a result that resists seven
falsifications as strengthened, not weakened, interpret the residual overlap
on unknown facts as the visual signature of unresolved superposition, and
regard the axis as a light, honest, almost-free structural probe to sit
alongside generation. Natural next steps, partly realized in the subsequent
work summarized in Sec.~\ref{sec:generalization}, are the multi-model study
of the dimensionality law, a knowledge-controlled dataset built from a
behavioral label, and the category-level structure of the axis, each preceded
by a behavioral feasibility check of the kind described in
Sec.~\ref{sec:future}.

\section*{Acknowledgments}

The author thanks Boaz Nadler for detailed and constructive comments on an
earlier version of this manuscript, which prompted, among many improvements,
the formalization of Section~\ref{sec:method} and Lemma~\ref{lem:labels}.

\begin{thebibliography}{99}

\bibitem{burns2023} C.~Burns, H.~Ye, D.~Klein, J.~Steinhardt.
\emph{Discovering Latent Knowledge in Language Models Without Supervision.}
arXiv:2212.03827, 2022 (ICLR 2023).

\bibitem{marks2024} S.~Marks, M.~Tegmark.
\emph{The Geometry of Truth: Emergent Linear Structure in Large Language Model Representations of True/False Datasets.}
arXiv:2310.06824, 2023 (COLM 2024).

\bibitem{li2023} K.~Li, O.~Patel, F.~Vi\'egas, H.~Pfister, M.~Wattenberg.
\emph{Inference-Time Intervention: Eliciting Truthful Answers from a Language Model.}
NeurIPS, 2023.

\bibitem{zou2023} A.~Zou et al.
\emph{Representation Engineering: A Top-Down Approach to AI Transparency.}
arXiv:2310.01405, 2023.

\bibitem{burger2024} L.~B\"urger, F.~A.~Hamprecht, B.~Nadler.
\emph{Truth is Universal: Robust Detection of Lies in LLMs.}
arXiv:2407.12831, 2024 (NeurIPS 2024).

\bibitem{cunningham2023} H.~Cunningham, A.~Ewart, L.~Riggs, R.~Huben, L.~Sharkey.
\emph{Sparse Autoencoders Find Highly Interpretable Features in Language Models.}
arXiv:2309.08600, 2023.

\bibitem{bricken2023} T.~Bricken et al.
\emph{Towards Monosemanticity: Decomposing Language Models With Dictionary Learning.}
Transformer Circuits Thread, Anthropic, 2023.

\bibitem{elhage2022} N.~Elhage et al.
\emph{Toy Models of Superposition.}
Transformer Circuits Thread, Anthropic, 2022 (arXiv:2209.10652).

\bibitem{barshalom2025} G.~Bar-Shalom, F.~Frasca, Y.~Galron, Y.~Ziser, H.~Maron.
\emph{Beyond Token Probes: Hallucination Detection via Activation Tensors with ACT-ViT.}
arXiv:2510.00296, 2025.

\bibitem{icr2025} \emph{ICR Probe: Tracking Hidden State Dynamics for Reliable Hallucination Detection in LLMs.}
arXiv:2507.16488, 2025.

\bibitem{hsad2025} \emph{LLM Hallucination Detection: A Fast Fourier Transform Method Based on Hidden Layer Temporal Signals (HSAD).}
arXiv:2509.13154, 2025.

\bibitem{counterfact} K.~Meng, D.~Bau, A.~Andonian, Y.~Belinkov.
\emph{Locating and Editing Factual Associations in GPT.}
NeurIPS, 2022. Dataset mirror: \texttt{NeelNanda/counterfact-tracing}.

\bibitem{truthfulqa} S.~Lin, J.~Hilton, O.~Evans.
\emph{TruthfulQA: Measuring How Models Mimic Human Falsehoods.}
ACL, 2022.

\end{thebibliography}

\end{document}
